% ============================================================================
% ZUSATZ-ARBEITSBLATT: Reibungskräfte (Rollreibung, Reifendruck, ABS)
% Physik Klasse 7 - Kräfte
% Stand: 05.02.2026
% ============================================================================

\documentclass[11pt,a4paper]{article}

\usepackage[utf8]{inputenc}
\usepackage[T1]{fontenc}
\usepackage[ngerman]{babel}
\usepackage{geometry}
\usepackage{helvet}
\usepackage{enumitem}
\usepackage{fancyhdr}
\usepackage{xcolor}
\usepackage{tcolorbox}
\usepackage{amssymb}
\usepackage{siunitx}
\usepackage{array}
\usepackage{qrcode}
\usepackage{textcomp}
\usepackage[hidelinks]{hyperref}

\sisetup{locale=DE, per-mode=power, inter-unit-product=\ensuremath{\cdot}}
\renewcommand{\familydefault}{\sfdefault}

\geometry{left=1.4cm, right=1.4cm, top=1.3cm, bottom=1.0cm}

% Fachfarbe Physik
\definecolor{physik}{HTML}{2c5aa0}
\colorlet{fachfarbe}{physik}
\definecolor{protokollrand}{HTML}{2a7a4b}
\definecolor{merkerand}{HTML}{b35c00}

% Kopfzeile
\pagestyle{fancy}
\fancyhf{}
\fancyhead[L]{\small\textbf{Physik Klasse 7} | Reibungskräfte -- Zusatz}
\fancyhead[R]{\small Zusatz-Arbeitsblatt}
\renewcommand{\headrulewidth}{0.4pt}
\setlength{\headheight}{14pt}

% Boxen
\tcbuselibrary{skins}

\newtcolorbox{merkkasten}[1][]{
    colback=white, colframe=fachfarbe,
    coltitle=black, colbacktitle=white,
    fonttitle=\bfseries\small,
    boxrule=1pt, arc=0pt,
    left=6pt, right=6pt, top=4pt, bottom=4pt,
    #1
}

\newtcolorbox{protokollbox}[1][]{
    colback=white, colframe=protokollrand,
    coltitle=black, colbacktitle=white,
    fonttitle=\bfseries\small,
    boxrule=1pt, arc=0pt,
    left=6pt, right=6pt, top=4pt, bottom=4pt,
    #1
}

\newtcolorbox{merkebox}[1][]{
    colback=white, colframe=merkerand,
    coltitle=black, colbacktitle=white,
    fonttitle=\bfseries\small,
    boxrule=1pt, arc=0pt,
    left=6pt, right=6pt, top=4pt, bottom=4pt,
    #1
}

\newcommand{\aufgabe}[1]{\noindent\textbf{\textcolor{fachfarbe}{#1}}}
\newcommand{\luecke}[1]{\underline{\hspace{#1}}}
\newcommand{\stern}[1]{\textcolor{merkerand}{\textbf{#1}}}

\setlength{\parindent}{0pt}
\setlength{\parskip}{0pt}

\begin{document}

% ============================================================================
% TITEL
% ============================================================================
\begin{minipage}[t]{0.80\textwidth}
    \vspace{0pt}
    {\Large\bfseries Reibungskräfte -- Zusatz \stern{$\star\star\star$}}\\[0.2em]
    {\small Name: \luecke{4.5cm} \hfill Datum: \luecke{2.5cm}}
\end{minipage}
\hfill
\begin{minipage}[t]{0.15\textwidth}
    \vspace{0pt}
    \raggedleft
    \qrcode[height=1.4cm]{https://devbox42.github.io/sciencesim/physik/kl07/02-kraefte/05-reibungskraefte/LP-05-komplett.html}
\end{minipage}
\vspace{0.1em}
\hrule
\vspace{0.4em}

% ============================================================================
% B1: ROLLREIBUNG
% ============================================================================
\begin{protokollbox}[title=B5: Rollreibung -- Warum Stahlräder bei Zügen?]
\small
\textbf{Beobachtung in der Simulation:} Je mehr sich ein Rad verformt (``walkt''), desto \luecke{2cm} ist die Rollreibung.

\vspace{0.3em}
\renewcommand{\arraystretch}{1.3}
\begin{tabular}{|p{3cm}|p{2.5cm}|p{2.5cm}|p{4cm}|}
\hline
\textbf{Material} & \textbf{Verformung} & \textbf{$F_{\text{Roll}}$} & \textbf{Einsatz} \\
\hline
Stahlrad & sehr gering & \luecke{1.5cm} N & \luecke{2.5cm} \\
\hline
Gummi (prall) & gering & \luecke{1.5cm} N & \luecke{2.5cm} \\
\hline
Gummi (platt) & groß & \luecke{1.5cm} N & -- \\
\hline
\end{tabular}

\vspace{0.3em}
\textbf{Erklärung:} Züge haben Stahlräder, weil \luecke{5.5cm}.

Autos haben Gummireifen, weil \luecke{5.5cm}.
\end{protokollbox}

\vspace{0.3em}

% ============================================================================
% B2: REIFENDRUCK
% ============================================================================
\begin{protokollbox}[title=B6: Reifendruck und Haftung]
\small
\textbf{Simulation: Verändere den Reifendruck und beobachte die Auswirkungen.}

\vspace{0.3em}
\renewcommand{\arraystretch}{1.3}
\begin{tabular}{|p{3cm}|p{2.5cm}|p{2.5cm}|p{3.5cm}|}
\hline
\textbf{Reifendruck} & \textbf{Kontaktfläche} & \textbf{Rollreibung} & \textbf{Haftung} \\
\hline
zu niedrig & \luecke{1.5cm} & \luecke{1.5cm} & \luecke{2cm} \\
\hline
optimal & mittel & niedrig & gut \\
\hline
zu hoch & \luecke{1.5cm} & \luecke{1.5cm} & \luecke{2cm} \\
\hline
\end{tabular}

\vspace{0.3em}
\textbf{Fazit:} Der optimale Reifendruck ist ein Kompromiss zwischen \luecke{3cm} und \luecke{3cm}.
\end{protokollbox}

\vspace{0.3em}

% ============================================================================
% B3: ABS
% ============================================================================
\begin{merkebox}[title=B7: ABS -- Das Antiblockiersystem \stern{$\star\star\star$}]
\small
\textbf{Problem ohne ABS:} Wenn die Räder \luecke{2.5cm}, wechselt der Reifen von \luecke{2.5cm} zu \luecke{2.5cm}. Da $F_{\text{Gleit}}$ \luecke{1cm} $F_{\text{Haft}}$ ist, wird der Bremsweg \luecke{2cm}!

\vspace{0.3em}
\textbf{Lösung mit ABS:} ABS erkennt das Blockieren und \luecke{3.5cm}, sodass das Rad wieder \luecke{2.5cm} und die \luecke{3cm} wirkt.

\vspace{0.3em}
\textbf{Ergebnis:} Mit ABS ist der Bremsweg \luecke{2cm} als ohne ABS, besonders auf \luecke{2.5cm} Straßen.
\end{merkebox}

\vspace{0.3em}

% ============================================================================
% T1: ABS MESSWERTE
% ============================================================================
\aufgabe{T3: Messwerte aus der ABS-Simulation (bei 50 km/h)}

\vspace{0.2em}
{\small
\renewcommand{\arraystretch}{1.3}
\begin{tabular}{|p{3cm}|p{3.5cm}|p{3.5cm}|p{3.5cm}|}
\hline
\textbf{Straße} & \textbf{Ohne ABS} & \textbf{Mit ABS} & \textbf{Unterschied} \\
\hline
Trocken & \luecke{2cm} m & \luecke{2cm} m & \luecke{2cm} m \\
\hline
Nass & \luecke{2cm} m & \luecke{2cm} m & \luecke{2cm} m \\
\hline
Vereist & \luecke{2cm} m & \luecke{2cm} m & \luecke{2cm} m \\
\hline
\end{tabular}
}

\vspace{0.2em}
\textbf{Erkenntnis:} Der Unterschied ist am größten auf \luecke{3cm} Straßen.

\vspace{0.4em}

% ============================================================================
% SELBSTCHECK
% ============================================================================
\begin{merkkasten}
\textbf{Selbstcheck: Das kann ich jetzt!} \hfill {\scriptsize \textcircled{\scriptsize 1} = unsicher \quad \textcircled{\scriptsize 2} = geht so \quad \textcircled{\scriptsize 3} = sicher}

\vspace{0.2em}
\small
\renewcommand{\arraystretch}{1.1}
\begin{tabular}{p{10.5cm}ccc}
Ich kann erklären, warum Züge Stahlräder haben. & $\square$ & $\square$ & $\square$ \\
Ich weiß, wie der Reifendruck die Haftung beeinflusst. & $\square$ & $\square$ & $\square$ \\
\stern{$\star\star\star$} Ich kann erklären, wie ABS funktioniert. & $\square$ & $\square$ & $\square$ \\
\end{tabular}
\end{merkkasten}

\end{document}
